% !TEX root =  master.tex
%      HYPERREF

%%%%%%%%%%%%%%%%%%%%%%%%%%%%%%%%%%%%%%%%%%%%%%%%%%%%%%%%%%
%	ANLEITUNG: 
%
% Passen Sie alle Stellen im Dokument an, die mit 
% @stud markiert sind
%
%%%%%%%%%%%%%%%%%%%%%%%%%%%%%%%%%%%%%%%%%%%%%%%%%%%%%%%%%%

\usepackage{makeidx}         % allows index generation
\usepackage{listings}	%Format Listings properly
\usepackage{lipsum}    %Blindtext
\usepackage{graphicx} % use various graphics formats
\usepackage[german]{varioref} 	% nicer references \vref
\usepackage{caption}	%better Captions
\usepackage{booktabs} %nicer Tabs
\usepackage{array}
\usepackage{chngcntr}
\usepackage[hidelinks=true]{hyperref} % keine roten Markierungen bei Links
\usepackage{fnpct} % Correct superscripts 
\usepackage[T1]{fontenc}
\usepackage[utf8]{inputenc}
\usepackage{calc} % Used for extra space below footsepline
\usepackage{acronym}
\usepackage{algorithm}
\usepackage{algpseudocode}
\usepackage{setspace}
\usepackage{wrapfig}
\usepackage{tabularx}

%
% @stud
%
%	FONT SELECTION: Entweder 1) Latin Modern oder 2) Times / Helvetica
\usepackage{lmodern}             % 1) Latin modern font
%\usepackage{mathptmx}           % 2) Helvetica / Times New Roman fonts (2 lines)
%\usepackage[scaled=.92]{helvet} % 2) Helvetica / Times New Roman fonts (2 lines)

%
% @stud
%
%	LANGUAGE SETTINGS
\usepackage[ngerman]{babel} 	        % german language
\usepackage[german=quotes]{csquotes} 	% correct quoting using \enquote{}
%\usepackage[english]{babel}          % english language
%\usepackage{csquotes} 	              % correct quoting using \enquote{}

%
% @stud
%
% Uncomment the following lines to support hard URL breaks in bibliography 
%\apptocmd{\UrlBreaks}{\do\f\do\m}{}{}
%\setcounter{biburllcpenalty}{9000}% Kleinbuchstaben
%\setcounter{biburlucpenalty}{9000}% Großbuchstaben

%
% @stud
%
%	FOOTNOTES: Count footnotes over chapters
%1 \counterwithout{footnote}{chapter}

%	ACRONYMS
\makeatletter
\@ifpackagelater{acronym}{2015/03/20}
{\renewcommand*{\aclabelfont}[1]{\textbf{{\acsfont{#1}}}}}{}
\makeatother

%	LISTINGS
\renewcommand{\lstlistingname}{Quelltext} 
\renewcommand{\lstlistlistingname}{Quelltextverzeichnis}
\lstset{numbers=left,
	numberstyle=\tiny,
	captionpos=b,
	basicstyle=\ttfamily\small}

%	ALGORITHMS
\renewcommand{\listalgorithmname}{Algorithmenverzeichnis }
\floatname{algorithm}{Algorithmus}

%		PAGE HEADER / FOOTER
%	    Warning: There are some redefinitions throughout the master.tex-file!  DON'T CHANGE THESE REDEFINITIONS!
\RequirePackage[automark]{scrlayer-scrpage}
%alternatively with separation lines: \RequirePackage[automark,headsepline,footsepline]{scrlayer-scrpage}

%\renewcommand*{\pnumfont}{\upshape\sffamily}
%\renewcommand*{\headfont}{\upshape\sffamily}
%\renewcommand*{\footfont}{\upshape\sffamily}

\renewcommand{\chaptermarkformat}{}
\RedeclareSectionCommand[beforeskip=0pt]{chapter}
\clearscrheadfoot

%\ifoot[\rule{0pt}{\ht\strutbox+\dp\strutbox}DHBW Mannheim]{\rule{0pt}{\ht\strutbox+\dp\strutbox}DHBW Mannheim}
\ofoot[\rule{0pt}{\ht\strutbox+\dp\strutbox}\pagemark]{\rule{0pt}{\ht\strutbox+\dp\strutbox}\pagemark}
\ohead{\headmark}

\newcommand{\TitelDerArbeit}[1]{\def\DerTitelDerArbeit{#1}\hypersetup{pdftitle={#1}}}
\newcommand{\AutorDerArbeit}[1]{\def\DerAutorDerArbeit{#1}\hypersetup{pdfauthor={#1}}}
\newcommand{\Firma}[1]{\def\DerNameDerFirma{#1}}
\newcommand{\Kurs}[1]{\def\DieKursbezeichnung{#1}}
\newcommand{\Abteilung}[1]{\def\DerNameDerAbteilung{#1}}
\newcommand{\Studiengangsleiter}[1]{\def\DerStudiengangsleiter{#1}}
\newcommand{\WissBetreuer}[1]{\def\DerWissBetreuer{#1}}
\newcommand{\FirmenBetreuer}[1]{\def\DerFirmenBetreuer{#1}}
\newcommand{\Bearbeitungszeitraum}[1]{\def\DerBearbeitungszeitraum{#1}}
\newcommand{\Abgabedatum}[1]{\def\DasAbgabedatum{#1}}
\newcommand{\Matrikelnummer}[1]{\def\DieMatrikelnummer{#1}}
\newcommand{\Studienrichtung}[1]{\def\DieStudienrichtung{#1}}
\newcommand{\ArtDerArbeit}[1]{\def\DieArtDerArbeit{#1}}
\newcommand{\Literaturverzeichnis}{Literaturverzeichnis}

\usepackage[backend=biber, autocite=footnote, style=authoryear, dashed=false]{biblatex}
\DefineBibliographyStrings{ngerman}{  %Change u.a. to et al. (german only!)
	andothers = {{et\,al\adddot}},
}
\setlength{\bibparsep}{\parskip}		%add some space between biblatex entries in the bibliography
\addbibresource{bibliography.bib}	%Add file bibliography.bib as biblatex resource
\usepackage{chngcntr}
\counterwithout{footnote}{chapter}


\newcommand{\setTitlepage}{
	% !TEX root =  master.tex
\begin{titlepage}
	\begin{center}
		\includegraphics[width=8cm]{\imagedir/logo.jpg}	
	\end{center}
	\vspace{2em}
	%\sffamily
	\begin{center}
		{\textsf{\large Duale Hochschule Baden-W\"urttemberg Mannheim}}\\[4em]
		{\textsf{\textbf{\large{\DieArtDerArbeit}arbeit}}}\\[6mm]
		{\textsf{\textbf{\Large{}\DerTitelDerArbeit}}} \\[1.5cm]
		{\textsf{\textbf{\large{}Studiengang Wirtschaftsinformatik}}\\[6mm]
		\textsf{\textbf{Studienrichtung \DieStudienrichtung}}}\\[6mm]
		\textsf{Kurs: \DieKursbezeichnung} \\[3cm]
		\begin{center}
			\begin{table}[h]
				\centering
				\begin{tabular}{c}
					Aaron Schweig \\
					Troy Kessler  \\
					Matthias Vonend  \\
					Michael Angermeier \\
					Patrick Mischka \\
					Jan Grübener \\
				\end{tabular}
			\end{table}
		\end{center}

		% \begin{minipage}{\textwidth}
		% 	\begin{tabbing}
		% 		Bearbeitungszeitraum: \hspace{0.85cm}\=\kill
		% 		% TODO: Verfasser
		% 		Kurs: \> \DieKursbezeichnung \\[1.5mm]
		% 		Bearbeitungszeitraum: \> \DerBearbeitungszeitraum\\[1.5mm]
		% %		alternativ:\\[1.5mm]
		% %		Eingereicht: \> \DasAbgabedatum	
		% 	\end{tabbing}
		% \end{minipage}
	\end{center}
\end{titlepage}
	\pagenumbering{roman} % Römische Seitennummerierung
	\normalfont	
}

%
% @stud
%
\newcommand{\settingLists}{
	%	Inhaltsverzeichnis
	\tableofcontents
	%	Abbildungsverzeichnis
	\listoffigures
	%	Tabellenverzeichnis
	\listoftables
	%	Listingsverzeichnis / Quelltextverzeichnis
	\lstlistoflistings
	% Algorithmenverzeichnis
	% \listofalgorithms
}

\newcommand{\initializeText}{
	\clearpage
	\ihead{\chaptername~\thechapter} % Neue Header-Definition
	\pagenumbering{arabic}           % Arabische Seitenzahlen
}

\newcommand{\initializeBibliography}{
	\ihead{}
	\printbibliography[title=\Literaturverzeichnis] 
	\cleardoublepage
}

\newcommand{\initializeAppendix}{
	\appendix
	\ihead{\appendixname~\thechapter}
}

\newcommand{\setword}[2]{%
  \phantomsection
  #1\def\@currentlabel{\unexpanded{#1}}\label{#2}%
}


\usepackage{color}

\definecolor{mygreen}{rgb}{0,0.6,0}
\definecolor{mygray}{rgb}{0.5,0.5,0.5}
\definecolor{mymauve}{rgb}{0.58,0,0.82}

\lstset{ 
  backgroundcolor=\color{white},
  breakatwhitespace=false,
  breaklines=true,
  captionpos=b,
  commentstyle=\color{mygreen},
  extendedchars=true,
  frame=single,
  keepspaces=false,
  keywordstyle=\color{blue},
  language=Java,
  numbers=left,
  numbersep=5pt,
  numberstyle=\color{mygray},
  rulecolor=\color{black},
  showspaces=false,
  showstringspaces=false,
  showtabs=false,
  stepnumber=2,
  stringstyle=\color{mymauve},
  tabsize=2,
  title=\lstname,
  basicstyle=\scriptsize
}

\usepackage{xspace}
% \usepackage{pifont}

\usepackage{tikz}
\newcounter{barcount}
\newenvironment{barchart}[1]{ % The only parameter is the maximum bar width, in cm
	\newcommand{\barwidth}{0.35}
	\newcommand{\barsep}{0.55}

	% Command to add a bar to the bar chart
	\newcommand{\baritem}[2]{ % The first argument is the bar label and the second is the percentage the current bar should take up of the total width
		\pgfmathparse{##2}
		\let\perc\pgfmathresult

		\pgfmathparse{#1}
		\let\barsize\pgfmathresult

		\pgfmathparse{\barsize*##2/5}
		\let\barone\pgfmathresult

		\pgfmathparse{(\barwidth*\thebarcount)+(\barsep*\thebarcount)}
		\let\barx\pgfmathresult

		\filldraw[fill=none, draw=black!] (0,-\barx) rectangle (5,-\barx-\barwidth);
		\filldraw[fill=black!90, draw=black!90] (0,-\barx) rectangle (\barone,-\barx-\barwidth);

		\node [label=180:{\textcolor{black}{##1: ##2/5}}] at (0,-\barx-0.175) {};
		\addtocounter{barcount}{1}
	}

	\newcommand{\baritemNL}[2]{ % The first argument is the bar label and the second is the percentage the current bar should take up of the total width
		\pgfmathparse{##2}
		\let\perc\pgfmathresult

		\pgfmathparse{#1}
		\let\barsize\pgfmathresult

		\pgfmathparse{\barsize*##2/5}
		\let\barone\pgfmathresult

		\pgfmathparse{(\barwidth*\thebarcount)+(\barsep*\thebarcount)}
		\let\barx\pgfmathresult

		\filldraw[fill=none, draw=black!] (0,-\barx) rectangle (5,-\barx-\barwidth);
		\filldraw[fill=black!90, draw=black!90] (0,-\barx) rectangle (\barone,-\barx-\barwidth);

		\node [label=180:{}] at (0,-\barx) {};
		\addtocounter{barcount}{1}
	}

	\begin{tikzpicture}
	\setcounter{barcount}{0}
}{
	\end{tikzpicture}
	\vspace{0.3cm}
}
